% Tailored 3-page AT&T Principal Data/AI Engineering resume; sourced from verified portfolio experience.
% Target: AT&T Principal Data/AI Engineering, R-123056-1 (Middletown NJ / Plano TX).
\documentclass[9pt,letterpaper]{extarticle}
\usepackage[left=0.68in,right=0.68in,top=0.45in,bottom=0.45in]{geometry}
\usepackage[T1]{fontenc}
\input{glyphtounicode}
\pdfgentounicode=1
\usepackage[utf8]{inputenc}
\usepackage{lmodern}
\usepackage{microtype}
\usepackage{pifont}
\usepackage{xcolor}
\usepackage{enumitem}
\usepackage{titlesec}
\usepackage[hidelinks]{hyperref}
\usepackage{parskip,ragged2e,needspace}
\pagestyle{empty}
\setlength{\parindent}{0pt}
\setlength{\parskip}{1pt}
\definecolor{accent}{HTML}{155B8A}
\titleformat{\section}{\large\bfseries\color{accent}}{}{0pt}{}[\vspace{0.5pt}\color{accent}\titlerule]
\titlespacing*{\section}{0pt}{3pt}{2pt}
\setlist[itemize]{leftmargin=1.15em,itemsep=0.7pt,topsep=1pt,parsep=0pt,partopsep=0pt}
\newcommand{\role}[4]{\Needspace{4\baselineskip}\textbf{#1}\hfill\textbf{#2}\\[-1pt]\textit{#3}\hfill\textit{#4}\par}
\newcommand{\project}[3]{\Needspace{5\baselineskip}\textbf{#1}\hfill\textit{#2 \textbar{} #3}\par}
\newcommand{\cert}[3]{\Needspace{4\baselineskip}\textbf{#1}\hfill\textit{#2}\par{\small #3\par}}
\hypersetup{
  pdftitle={Waseem Raza, PhD - AT&T Principal Data/AI Engineering R-123056-1},
  pdfauthor={Waseem Raza}
}


\begin{document}
\begin{center}
{\LARGE\bfseries Waseem Raza, PhD}\\[2pt]
{\large\bfseries Principal Data/AI Engineering | AI Platforms, Data Pipelines \& MLOps}\\[1pt]
{\small Production-Oriented AI | Azure Data Engineering | GenAI Workflows | Model Reliability}\\[3pt]
{\footnotesize New Jersey, United States \quad
\href{tel:+17328039703}{+1 732-803-9703} \quad
\href{mailto:waseem.placements@gmail.com}{waseem.placements@gmail.com}\\[1pt]
\href{https://www.linkedin.com/in/waseem-raza-phd-9a737858/}{LinkedIn} \textbar{}
\href{https://github.com/WaseemRaza844}{GitHub} \textbar{}
\href{https://waseemraza844.github.io/portfolio/}{Portfolio}}
\end{center}
\vspace{-5pt}

\section{Professional Summary}
AI/ML and data engineering professional with a PhD in Electrical and Computer Engineering and \textbf{7+ years of combined engineering, applied research and analytics experience}. Build, validate and operationally integrate data-driven solutions using Python, SQL, Azure Databricks, Azure Functions and Synapse Analytics. At AT\&T, develop AI-based E911 diagnostics, improve root-cause classification accuracy from \textbf{20\% to 80\%}, and evaluate performance across heterogeneous operational domains. Previously reduced cloud ETL preparation time by \textbf{40\%} and improved analytical-output reliability by \textbf{20\%}. Strengths include scalable ingestion, structured/unstructured data, cross-domain model evaluation, incident-oriented troubleshooting, GenAI prototypes and collaboration with software, network and business stakeholders. Seeking to translate these skills into enterprise AI/data platform engineering and operations.

\section{Technical Competencies}
\textbf{Data Engineering \& Integration:} Python, SQL, ETL/ELT, ingestion, cleansing, schema transformation, data profiling, feature engineering, validation checks, structured/semi-structured/unstructured data, AI-ready datasets.\\
\textbf{Cloud Data \& AI:} Hands-on Azure Databricks, Apache Spark/PySpark, Azure Functions, Synapse Analytics, Azure data pipelines and cloud analytics; AWS Bedrock and Google Cloud through training/projects.\\
\textbf{MLOps \& Model Reliability:} Train/validation/test workflows, model-performance evaluation, monitoring-oriented metrics, model drift/domain-shift testing, error analysis, reproducibility, failure-mode review and deployment readiness.\\
\textbf{AI, GenAI \& Workflows:} PyTorch, TensorFlow, scikit-learn, XGBoost, NLP, LLM-assisted extraction, RAG, embeddings/vector retrieval, LangChain, LangGraph and agent orchestration (portfolio / coursework as noted).\\
\textbf{Platform Engineering Tools:} Git/GitHub, GitHub Actions, CI/CD, Docker, Linux, REST APIs, Flask, Jupyter, test automation, debugging, logging/data-quality checks and system-design fundamentals; Kubernetes/OpenShift familiarity.\\
\textbf{Analytics \& Communication:} Pandas, NumPy, SciPy, Matplotlib, Plotly; stakeholder requirements, technical documentation, workflow troubleshooting, performance reporting and cross-functional delivery.

\section{Professional Experience}
\role{Machine Learning Engineer}{02/2025 -- Present}{AT\&T (contract assignment via Innovatix Inc.)}{Bedminster, NJ}
\begin{itemize}
\item Develop, test and integrate AI-driven diagnostics for critical E911 service-assurance operations; improved automated root-cause classification accuracy from \textbf{20\% to 80\%}.
\item Design reusable preprocessing and cross-domain ML evaluation workflows for text, location, network context and operational signals; strengthen reliability under distribution shifts.
\item Develop multimodal GenAI-assisted prototypes for VoNR/5G diagnostic use cases, combining structured context and unstructured work-order information.
\item Perform systematic failure analysis, data validation and performance benchmarking to diagnose model weaknesses and support production-oriented integration decisions.
\item Collaborate with network engineers, software teams and public-safety stakeholders to translate operational problems into validated datasets, measurable AI objectives and practical engineering workflows.
\end{itemize}

\role{Data Scientist / AI-ML Engineer}{09/2024 -- 12/2024}{Impac Exploration Services}{Weatherford, OK}
\begin{itemize}
\item Designed Azure-based ingestion, transformation and data-quality workflows for large geochemical datasets using \textbf{Python, Databricks, Azure Functions and Synapse Analytics}.
\item Automated repeatable ETL/data-preparation processes, reducing preparation time by \textbf{40\%}; introduced cross-validation and consistency checks that increased analytical-output reliability by \textbf{20\%}.
\item Built LLM-assisted structured extraction and analytics pipelines for scientific reports and well data, translating domain requirements into AI-ready datasets.
\item Worked directly with geologists to troubleshoot data quality issues, review outputs and turn research questions into usable analytics pipelines.
\end{itemize}

\section{Additional Engineering \& Applied Research Experience}
\role{Advanced-Degrees Research Intern}{06/2024 -- 08/2024}{AT\&T Labs}{Austin, TX}
\begin{itemize}
\item Analyzed \textbf{500K+ LTE/5G diagnostic logs} to investigate anomalies and network behavior; applied ML methods that reduced diagnostic analysis time by \textbf{20\%}.
\item Built quantitative evaluation and reporting workflows for large operational datasets and communicated performance findings to engineering teams.
\end{itemize}

\role{Graduate Research Assistant -- AI/ML Systems}{09/2019 -- 07/2025}{AI4Networks, University of Oklahoma}{Norman, OK}
\begin{itemize}
\item Designed reusable Python simulation, data-generation and experiment pipelines for ML development, including 3GPP-aligned network datasets and scarce-data augmentation.
\item Implemented GAN, autoencoder, transfer/domain-adaptation and reinforcement-learning workflows; conducted \textbf{100+ controlled experiments} with systematic model comparison.
\item Built robustness-testing pipelines for changing data distributions and missing labels, combining repeatable evaluation metrics, quantitative error analysis and technical reporting.
\item Mentored research collaborators on reproducibility, experiment tracking, debugging, code reviews and publication-quality results.
\end{itemize}

\role{Advanced Wireless Technology Development Intern}{09/2022 -- 01/2023}{Futurewei Technologies}{Schaumburg, IL}
\begin{itemize}
\item Improved autoencoder-based wireless channel-data compression by \textbf{25\%} using transfer learning; completed \textbf{50+ benchmark experiments} across operating conditions.
\item Developed repeatable feature preparation and performance-comparison procedures for high-dimensional datasets and translated findings into engineering recommendations.
\end{itemize}

\section{Selected Enterprise AI \& Data Platform Projects}
\project{AI-Ready Lakehouse and Data Validation Pipeline}{Portfolio / Applied Project}{2024}
Designed ingestion, schema transformation, cleansing, validation and feature-preparation workflows for structured and semi-structured data. Built scalable processing patterns in Azure Databricks and Spark/PySpark to prepare reliable datasets for downstream analytics and ML.\\[-1pt]
{\small\textit{Stack:} Python, SQL, Azure Databricks, Spark/PySpark, ETL/ELT, data-quality validation, cloud data engineering.}

\project{GenAI-Assisted Operational Diagnostics}{Industry Initiative}{2025--Present}
Combined work-order text, location and network context to prototype GenAI-assisted diagnostics alongside cross-domain predictive models. Focused on traceable preprocessing, evaluation against operating-condition changes, failure-case analysis and collaboration with engineering users.\\[-1pt]
{\small\textit{Stack:} Python, NLP, structured/unstructured data, multimodal features, ML evaluation, operational workflow integration.}

\project{Enterprise Document and Unstructured-Data Processing}{Applied / Guided Project}{2025--2026}
Designed document-ingestion, text extraction, chunking, normalization and structured-output workflows, connecting unstructured content to downstream analysis and retrieval. Evaluated data readiness, extraction consistency and reproducible output formats.\\[-1pt]
{\small\textit{Stack:} Python, OCR/document intelligence, LLM extraction, JSON, data validation, RAG, embeddings, vector retrieval.}

\project{Retrieval-Augmented and Agentic AI Workflows}{Portfolio / Certification Projects}{2026}
Implemented grounded RAG and tool-using agent workflows with document retrieval, embeddings, stateful routing and structured outputs. Explored how evaluation, human oversight and failure analysis support more dependable GenAI workflows.\\[-1pt]
{\small\textit{Stack:} LangChain, LangGraph, LlamaIndex, ChromaDB, FAISS, AWS Bedrock, model/tool integration.}

\project{SyntheticNET: Reproducible Large-Scale Data Simulation}{University Research Platform}{2020--2025}
Created repeatable Python-based data generation and experimental pipelines for network-state, propagation and mobility scenarios. Enabled controlled datasets, model benchmarking and reliability tests where ground-truth measurements were limited.\\[-1pt]
{\small\textit{Stack:} Python, numerical simulation, repeatable pipelines, synthetic data, robust ML evaluation.}

\section{Production Engineering, Operational Excellence \& Collaboration}
\begin{itemize}
\item \textbf{Reliability:} Review data lineage and input consistency, test cross-domain model behavior, inspect classification errors and document failure patterns before operational changes.
\item \textbf{Automation:} Use reusable Python processing stages, source control, automated checks and CI/CD concepts to reduce manual analysis and support repeatable deliveries.
\item \textbf{Data quality:} Apply schema/consistency checks, exception analysis and validation gates when moving structured or document-derived data toward ML and analytics.
\item \textbf{Collaboration:} Work across engineering, data science, software and domain teams; communicate tradeoffs, root causes, requirements and acceptance criteria in operational terms.
\end{itemize}

\section{Selected Certifications \& Professional Development}
\begin{itemize}
\item \textbf{Generative AI Applications with Amazon Bedrock} -- Coursera, completed 08/2026; foundation models, managed GenAI services and RAG workflows.
\item \textbf{Agentic AI with LangChain and LangGraph} -- Coursera, completed 02/2026; stateful agents, tool use, routing and orchestration.
\item \textbf{Generative AI Engineering and Fine-Tuning Transformers} -- Coursera, completed 08/2026; transformers, adaptation and model evaluation.
\item \textbf{Machine Learning} -- Stanford Online / Coursera, completed 03/2020; applied modeling, validation and optimization.
\end{itemize}

\section{Education}
\textbf{PhD, Electrical and Computer Engineering}, University of Oklahoma, Norman, OK \hfill 2025\\
\textbf{MSc, Telecommunication Engineering}, University of Engineering and Technology, Taxila \hfill 2016\\
\textbf{BSc, Telecommunication Engineering}, University of Engineering and Technology, Taxila \hfill 2014

\section{Selected Technical Research}
\textbf{AI-Powered Resilience: A Dual-Approach for Outage Management in Dense Cellular Networks} -- \textit{Computer Communications}, 2025.\\
\textbf{An AI-Driven Framework for Enhancing Resilience in Propagation Models to Enable Digital Twin} -- \textit{IEEE PIMRC}, 2024.\\
\textbf{On Multi-Parameter Optimization and Proactive Reliability in 5G and Beyond Cellular Networks} -- \textit{Sensors}, 2025.

\end{document}
